\documentclass{article}
\usepackage{graphicx} % Required for inserting images
\usepackage{amsmath}
\usepackage{tikz}

\title{Solutions to Worksheet \#4}
\author{}
\date{August 5, 2026}

\begin{document}

\maketitle
\begin{enumerate}
\item \textbf{Solution.}
Recall that
\[
X_G=\sum_{\kappa}\prod_{v\in V(G)}x_{\kappa(v)},
\]
where the sum is over proper colorings by positive integers.  Let
$\sigma$ be any permutation of the positive integers.  Replacing every
color $i$ by $\sigma(i)$ is a bijection from proper colorings to proper
colorings.  It permutes the variables in every monomial but leaves the
whole sum unchanged.  Thus $X_G$ is invariant under every permutation
of the variables, so it is symmetric.

\item \textbf{Solution.}
Let $e$ be any edge of $C_4$.  Deleting $e$ produces the path $P_4$,
which is a tree on four vertices, while contracting $e$ produces the
triangle $C_3$.  Therefore
\[
\chi_{C_4}(n)
=n(n-1)^3-n(n-1)(n-2).
\]
Simplifying,
\[
\boxed{\chi_{C_4}(n)
=n(n-1)(n^2-3n+3)
=n^4-4n^3+6n^2-3n.}
\]

\item \textbf{Solution.}
We use Stanley's power-sum formula
\[
X_G=\sum_{A\subseteq E(G)}(-1)^{|A|}p_{\lambda(A)},
\]
where $\lambda(A)$ is the partition of the component sizes of the
spanning subgraph $(V(G),A)$.
\begin{enumerate}
\item Deleting the chosen edge gives $P_4$.  Among its three edges,
there are three one-edge subsets, two adjacent two-edge subsets, one
disjoint two-edge subset, and one three-edge subset.  Hence
\[
\boxed{X_{P_4}
=p_{1,1,1,1}-3p_{2,1,1}+2p_{3,1}+p_{2,2}-p_4.}
\]

\item Contracting the edge gives $C_3=K_3$.  The subset formula gives
\[
\boxed{X_{C_3}=p_{1,1,1}-3p_{2,1}+2p_3.}
\]

\item The two terms in a deletion--contraction formula do not even
have the same degree: $X_{P_4}$ has degree $4$, while $X_{C_3}$ has
degree $3$.  More explicitly,
\[
X_{C_4}
=p_{1,1,1,1}-4p_{2,1,1}+4p_{3,1}+2p_{2,2}-3p_4,
\]
which is not $X_{P_4}-X_{C_3}$.  Thus the ordinary
deletion--contraction relation does not hold for the chromatic
symmetric function.
\end{enumerate}

\item \textbf{Solution.}
\begin{enumerate}
\item By the binomial theorem,
\[
\begin{aligned}
n(n-1)^{k-1}
&=n\sum_{i=0}^{k-1}\binom{k-1}{i}n^i(-1)^{k-1-i}\\
&=\boxed{\sum_{j=1}^{k}(-1)^{k-j}
\binom{k-1}{j-1}n^j}.
\end{aligned}
\]

\item Taking $j=k-1$, the coefficient is
\[
\boxed{-\binom{k-1}{k-2}=-(k-1).}
\]

\item For a simple graph,
\[
\boxed{[n^{k-1}]\chi_G(n)=-|E(G)|.}
\]
We prove this by induction on the number of edges.  For the graph with
no edges, $\chi_G(n)=n^k$, and both sides are $0$.  Choose an edge $e$.
By induction, the coefficient of $n^{k-1}$ in
$\chi_{G-e}(n)$ is $-(|E(G)|-1)$.  The contracted graph has $k-1$
vertices, so its chromatic polynomial is monic of degree $k-1$.
Consequently, deletion--contraction gives
\[
[n^{k-1}]\chi_G(n)
=-(|E(G)|-1)-1=-|E(G)|.
\]

\item If $\chi_G(n)=n(n-1)^{k-1}$, part (b) says that the coefficient
of $n^{k-1}$ is $-(k-1)$.  By part (c), $G$ has $k-1$ edges.  Since
$G$ is connected, a connected graph on $k$ vertices with $k-1$ edges
is a tree.
\end{enumerate}

\item \textbf{Solution.}
Call the graphs $G_1$ and $G_2$ in their displayed order.  In either
graph there are four independent pairs.  Thus a coloring of type
$(2,1,1,1)$ is obtained by choosing one of those pairs and then
assigning the three singleton colors in $3!$ ways, giving coefficient
$4\cdot3!=24$.

There are two ways in either graph to choose two disjoint independent
pairs.  Since the two pair-colors can be interchanged, the coefficient
of $m_{2,2,1}$ is $2\cdot2!=4$.  No independent set has size $3$, and
a coloring with all colors distinct contributes $5!=120$.  Therefore
the two graphs have the same answer:
\[
\boxed{X_{G_1}=X_{G_2}
=4m_{2,2,1}+24m_{2,1,1,1}+120m_{1,1,1,1,1}.}
\]

\item \textbf{Solution.}
In the power expansion, the coefficient of $p_n$ is
\[
\sum_{\substack{A\subseteq E(G)\\(V(G),A)\text{ connected}}}
(-1)^{|A|}.
\]
It is zero when $G$ is disconnected, because no edge subset can
connect all the vertices.  When $G$ is connected, the signed sum is
nonzero; up to the sign $(-1)^{n-1}$, it counts acyclic orientations
having a unique sink at any fixed vertex.  At least one such
orientation is obtained by ordering vertices by their distance from
the chosen sink and directing edges toward smaller distance.
Consequently,
\[
\boxed{G\text{ is connected if and only if }[p_n]X_G\ne0.}
\]
More generally, the smallest number of parts among partitions
indexing nonzero power-sum terms is the number of connected components.

\item \textbf{Solution.}
The term $p_{2,1^{n-2}}$ in the subset expansion can arise only from a
subset containing exactly one edge.  Every edge contributes
$-p_{2,1^{n-2}}$.  Hence
\[
\boxed{[p_{2,1^{n-2}}]X_G=-|E(G)|.}
\]
If $X_G=X_H$, these coefficients agree, so $G$ and $H$ have the same
number of edges.

\item \textbf{Solution.}
In a proper coloring of $K_n$, all vertices receive distinct colors.
For each set of $n$ colors there are $n!$ assignments to the vertices,
so
\[
\boxed{X_{K_n}=n!e_n.}
\]
In particular,
\[
X_{K_1}=e_1,\quad
X_{K_2}=2e_2,\quad
X_{K_3}=6e_3,\quad
X_{K_4}=24e_4,\quad
X_{K_5}=120e_5.
\]

\item \textbf{Solution.}
Let $e=uv$, where $u$ is the leaf and $v$ is its neighbor.
\begin{enumerate}
\item Deleting $e$ leaves $u$ isolated.  Thus
$G\setminus e=(G-u)\mathbin{\sqcup}K_1$.

\item Dot-contraction contracts $u$ and $v$ and then adds an isolated
vertex.  Contracting a leaf edge simply removes the leaf, so the
result is again $(G-u)\mathbin{\sqcup}K_1$.

\item The graphs in (a) and (b) are isomorphic.

\item Leaf-contraction contracts $u$ and $v$ and then attaches a new
leaf to the merged vertex.  This recreates a graph isomorphic to $G$.

\item The deletion-near-contraction relation is
\[
X_G=X_{G\setminus e}-X_{\text{dot}(G,e)}
    +X_{\text{leaf}(G,e)}.
\]
For a leaf edge, the first two terms cancel and the last term is
$X_G$, so the relation reduces to the tautology $X_G=X_G$.  It makes
no progress.  This is why the recursive algorithm uses internal
edges.
\end{enumerate}

\item \textbf{Solution.}
Write $\mathrm{st}_r=X_{St_r}$ and
$\mathrm{st}_\lambda=\prod_i\mathrm{st}_{\lambda_i}$.  There are
three trees on five vertices: $St_5$, the tree $T_5$ with degree
sequence $(3,2,1,1,1)$, and the path $P_5$.

For $T_5$, apply DNC to the unique internal edge adjacent to the
degree-$3$ vertex.  The deletion, dot-contraction, and
leaf-contraction terms are respectively star forests of types
$(3,2)$, $(4,1)$, and $(5)$.  Thus
\[
\boxed{X_{T_5}=\mathrm{st}_5-\mathrm{st}_{4,1}
                       +\mathrm{st}_{3,2}.}
\]
Similarly, applying DNC to the middle edge of $P_4$ gives
\[
X_{P_4}=\mathrm{st}_{2,2}-\mathrm{st}_{3,1}+\mathrm{st}_4.
\]
Applying DNC to an internal edge of $P_5$ then gives
\[
\begin{aligned}
X_{P_5}
&=\mathrm{st}_{3,2}-X_{P_4}\mathrm{st}_1+X_{T_5}\\
&=\boxed{\mathrm{st}_5-2\mathrm{st}_{4,1}
+2\mathrm{st}_{3,2}+\mathrm{st}_{3,1,1}
-\mathrm{st}_{2,2,1}.}
\end{aligned}
\]
Finally, by definition,
\[
\boxed{X_{St_5}=\mathrm{st}_5.}
\]

\item \textbf{Solution.}
Call the displayed graphs $H_1$ and $H_2$.  The subset formula gives
the same power-sum expansion for both:
\[
\begin{aligned}
X_{H_i}={}&p_{1^6}-6p_{2,1^4}+6p_{2,2,1,1}-p_{2,2,2}
+8p_{3,1,1,1}\\
&-6p_{3,2,1}-9p_{4,1,1}+3p_{4,2}
+6p_{5,1}-2p_6.
\end{aligned}
\]
To convert to the star basis, one may use
\[
p_n=\sum_{r=0}^{n-1}(-1)^r\binom{n-1}{r}
\mathrm{st}_{r+1,1^{n-1-r}}
\]
and multiply this identity for the parts of each power-sum index.
After collecting terms,
\[
\boxed{\begin{aligned}
X_{H_1}=X_{H_2}
={}&2\mathrm{st}_6-4\mathrm{st}_{5,1}
+3\mathrm{st}_{4,2}+2\mathrm{st}_{4,1,1}\\
&-3\mathrm{st}_{3,2,1}+\mathrm{st}_{2,2,2}.
\end{aligned}}
\]
Thus these two nonisomorphic graphs have the same chromatic symmetric
function.

\end{enumerate}
\end{document}
